Do some football clubs receive systematically more favorable refereeing decisions? I assemble match-level data on fouls, cards, penalties, and stoppage time for the five largest European leagues and the UEFA Champions League from 2001/02 to 2025/26, about 48,500 matches. I ask where FC Barcelona sits relative to every other club before and after 2018, the year in which the club's reported payments to the vice-president of Spain's referees' committee ended. Raw comparisons strongly favor Barcelona, but they largely reflect its possession-dominant style. Conditional on opponent, home status, pre-match odds, and possession, Barcelona receives 0.17--0.23 fewer yellow cards per match than the average club. Conditional on fouls called, it receives no fewer cards, and its penalty treatment is unremarkable: it ranks 93rd of 150 clubs on net penalties. Normalizing fouls by possession reverses the sign of its foul advantage. Barcelona's net card advantage in La Liga declines after 2017/18 relative to Real Madrid (by 0.57 cards per match), but the timing tracks its sporting decline and Messi's departure at least as closely as the end of the payments. The clearest favoritism is in stoppage time: referees add about 0.7 more minutes when Barcelona trails than when it leads. Real Madrid receives the same advantage, suggesting status-based social pressure rather than a club-specific arrangement.
